\documentclass[twocolumn]{tem}
\usepackage{zhlipsum}
% 页眉示例: 左/中/右三段, 位置由 tem.cls 里的制表位控制
\temSetHeaderContent{实验报告}{傻逼大学}{第\thepage 页}

% 如果你把字体放到 fonts/ 下, 文件名匹配即可自动生效.
% 例如: fonts/SimSun.ttf 和 fonts/Times New Roman.ttf

\begin{document}
\XtuTitleAuthorAddress{示例标题 Example Title}
{张三}
{(傻逼大学)}
{Example Title}
{Zhang San}
{Stupid University}

\Zhabstract{这里是中文摘要内容, 用于测试小五号宋体与 14pt 行距。}
\Zhkeyword{关键词1; 关键词2; 关键词3}
\Enabstract{This is an English abstract used to test 9pt Times New Roman with 14pt baseline skip.}
\Enkeyword{keyword1; keyword2; keyword3}


\section{示例标题}
\zhlipsum[1-4]


\subsection{代码}
\begin{lstlisting}[language=Python]
print('hello')
\end{lstlisting}

\subsection{算法}
\begin{algorithm}
\caption{示例算法}
\end{algorithm}

\end{document}
